%! Author = chtompki
%! Date = 1/14/26
\documentclass[10pt]{amsart}
% Default margins are too wide all the way around. I reset them here
\setlength{\hoffset}{-.75in}
\setlength{\textwidth}{6.5in}
\setlength{\voffset}{-.5in}
\setlength{\textheight}{9.0in}
\usepackage{amsmath}
\usepackage{amssymb}
\usepackage{amsthm}
\usepackage{lmodern}
\usepackage{microtype}
\usepackage{enumitem}
\usepackage{setspace}
\usepackage{graphicx}
\usepackage{hyperref}
\usepackage{tikz}

\newenvironment{solution}
{\begin{proof}[Solution]}
{\end{proof}}

\NewDocumentCommand{\vitem}{v}{%
    \item[\texttt{#1}]%
}

\title{Dynamical Systems\\Problem Set 1}
\author{Rob Tompkins}
\date{\today}

\begin{document}
    \maketitle
    \date{\today}
    \begin{onehalfspace}
        We begin by noting that all of the files referenced in this document are available at: \\
        \underline{https://github.com/chtompki/MathFall2026/tree/main/DynamicalSystems/ProblemSet1}

        \vspace{2.55em}
        {\bfseries\large Problem Set Instructions:}

        \begin{itemize}
            \item In your first attempt of the problem set problems, you are encouraged to treat the problem
            set as an open-notes quiz. Work through it on your own, without consulting other people,
            internet resources, or any solutions/answers. Work on each problem, completing as much as
            you are able to, and making a note in your work whenever you become stuck or confused.

            \item After your initial individual attempt, collaboration is encouraged as you continue to work on
            the problems. You can discuss the problem and potential approaches with your classmates,
            and use resources from class or other sources/references (including related journal articles).
            However, please do not look at the student solution manual or other posted solutions until
            after you have submitted your work in Gradescope.

            \item In your submission, please show all your work and include any code you create in addition the
            output of your computations. Be sure to acknowledge any substantial help you receive. This
            includes collaboration with others, references and online resources.
        \end{itemize}

        \vspace{2.4em}
        {\noindent\bfseries\large Part 1:}


            \noindent\textbf{1.} Become a student member of the Society for Industrial and Applied Mathematics (SIAM)
            (You don't need to turn anything in). Because VCU has a SIAM Student Chapter, you can
            become a member for free.

            \begin{enumerate}
                \item Find out more details and sign up at this link:\\
                \url{http://siam.org/membership/individual-membership/student-membership/}

                \item Look on the SIAM website for resources/events that might be of interest to you. Also
                take a look in the SIAM News section linked below for any interesting articles related to
                applications of dynamical systems:\\
                \url{https://www.siam.org/publications/siam-news/}
            \end{enumerate}

            \noindent\textbf{2.} Solving ODEs and phase portrait analysis with Matlab. You don't need to turn anything in
            for this, but use what you've learned to complete the computational aspects of the questions
            below. You can also feel free to use Python, Mathematica, or other software/programming
            languages instead of Matlab for your problem sets if you prefer. If you do want to use something
            different, please discuss this with me beforehand so that we can make sure you will still be
            able to do everything required over the course of the semester.

            \begin{enumerate}
                \item Get access to Matlab if you don't already have it. Information about downloading Matlab
                or using the online version as a VCU student can be found at the link below:\\
                \url{https://knowledgebase.vcu.edu/portal/app/portlets/results/viewsolution.jsp?solutionid=250303090503713}
            \end{enumerate}


          \begin{enumerate}[start=2]
            \item if you are unfamiliar with Matlab, go through the Matlab onramp tutorial linked below:\\
            \url{https://matlabacademy.mathworks.com/details/matlab-onramp/gettingstartedL}

            \item Complete the ``SlopeFields.mlx'' and ``Equilibria1D.mlx'' tutorials listed at the link below.
            In addition to reviewing some concepts related to one-dimensional dynamical systems,
            these tutorials provide example code for numerically analyzing 1D systems in Matlab.
            You may either download them or run them using Matlab Online.
            Link to Tutorials Main Page:
            \url{https://www.mathworks.com/matlabcentral/fileexchange/95513-qualitative-analysis-of-odes}

            \item Complete the following tutorial for using Matlab to implement the Explicit Euler Method
            for solving and initial value problem:\\
            \url{https://www.mathworks.com/videos/solving-odes-in-matlab-1-euler-ode1-117526.html}
          \end{enumerate}


          \noindent\textbf{3.} Computing bifurcation diagrams by numerical continuation in COCO. The numerical continuation
          package is written in Matlab, and will be useful for computing bifurcation diagrams.
          There are other options for those who prefer alternatives to Matlab, so please talk to me if
          that is the case.

          \begin{enumerate}
              \item download and install COCO using the instructions at the link below:\\
              \url{https://sourceforge.net/p/cocotools/wiki/Home/}

              \item Go through Sections 1--3.3 of the ``getting started'' tutorial, along with links to the re-
              lated Matlab tutorial file, which are both linked below (You don't need to turn anything
              in). Note that both the pdf and matlab file linked below are also included in your COCO
              installation directory.\\
              Tutorial: \url{https://drive.google.com/file/d/1I0Wm5p9UDiUEiFQ2IRPoUyvXF_N2w8D4/view?usp=sharing}\\
              Matlab Files: \url{https://drive.google.com/file/d/1p6eaDWMDErjFs7B1WBEj0OZrdskQk-kK/view?usp=drive_link}
          \end{enumerate}


        \vspace{0.45em}
        {\noindent\bfseries\large Part 2:}\quad from Strogatz, 2nd Edition, Chapter 2

        \noindent\textbf{2.2.4} Analyze the following equation graphically.
        Sketch the vector field on the real line find all the fixed points, classify their
        stability, and sketch the graph of $x(t)$ for different initial conditions.
        Then try to find the analytical solution for $x(t)$, or state that it's impossible.
        The equation in question is:
        \[
            \dot{x} = e^{-x}\sin(x)
        \]

        \begin{solution}
            Set
            \[
                f(x)=e^{-x}\sin x.
            \]
            Because $e^{-x}>0$ for every $x\in\mathbb R$, the sign of $f(x)$ is the
            same as the sign of $\sin x$. The fixed points are therefore
            \[
                f(x_*)=0
                \quad\Longleftrightarrow\quad
                \sin x_*=0
                \quad\Longleftrightarrow\quad
                \boxed{x_*=n\pi,\qquad n\in\mathbb Z.}
            \]
            Moreover,
            \[
                \dot{x}>0 \quad\text{on }(2k\pi,(2k+1)\pi),
                \qquad
                \dot{x}<0 \quad\text{on }((2k-1)\pi,2k\pi).
            \]
            Thus a representative portion of the phase line is
            \[
                \begin{tikzpicture}[baseline=-0.5ex,>=stealth,scale=0.95]
                  \draw[->] (-5.4,0)--(5.8,0) node[right] {$x$};
                  \foreach \x/\lab in {-4.5/$-2\pi$,-2.25/$-\pi$,0/$0$,2.25/$\pi$,4.5/$2\pi$}{
                      \fill (\x,0) circle (2pt);
                      \node[below=5pt] at (\x,0) {\lab};
                  }

                  % Chevrons placed directly on the phase line indicate the flow direction.
                  \node[fill=white,inner sep=0.5pt] at (-3.375,0) {$>$};
                  \node[fill=white,inner sep=0.5pt] at (-1.125,0) {$<$};
                  \node[fill=white,inner sep=0.5pt] at ( 1.125,0) {$>$};
                  \node[fill=white,inner sep=0.5pt] at ( 3.375,0) {$<$};


                \end{tikzpicture}
            \]

            For the stability classification, differentiate:
            \[
                f'(x)=e^{-x}(\cos x-\sin x).
            \]
            At a fixed point $x_*=n\pi$,
            \[
                f'(n\pi)=(-1)^n e^{-n\pi}.
            \]
            Consequently,
            \[
                \boxed{x=2k\pi\text{ is unstable},\qquad
                x=(2k+1)\pi\text{ is asymptotically stable}.}
            \]
            All the fixed points are hyperbolic.

            Now let $x(0)=x_0$. If $x_0=n\pi$, uniqueness gives the equilibrium
            solution
            \[
                x(t)\equiv n\pi.
            \]
            If $2k\pi<x_0<(2k+1)\pi$, then $x(t)$ increases monotonically and
            \[
                \lim_{t\to-\infty}x(t)=2k\pi,
                \qquad
                \lim_{t\to\infty}x(t)=(2k+1)\pi.
            \]
            If $(2k-1)\pi<x_0<2k\pi$, then $x(t)$ decreases monotonically and
            \[
                \lim_{t\to-\infty}x(t)=2k\pi,
                \qquad
                \lim_{t\to\infty}x(t)=(2k-1)\pi.
            \]
            Thus every nonconstant orbit connects an unstable even multiple of $\pi$
            to one of the adjacent stable odd multiples. Representative solution curves
            near $x=0$ are shown below.
            \[
                \begin{tikzpicture}[>=stealth,x=0.85cm,y=0.8cm]
                  \draw[->] (-4.8,0)--(5.2,0) node[right] {$t$};
                  \draw[->] (0,-4.25)--(0,4.25) node[above] {$x(t)$};
                  % Stable equilibrium solutions are solid; the unstable one is dotted.
                  \draw[thick] (-4.7,3.14)--(5.0,3.14)
                  node[right,align=left] {$\pi$\\[-1mm]{\scriptsize stable}};
                  \draw[thick] (-4.7,-3.14)--(5.0,-3.14)
                  node[right,align=left] {$-\pi$\\[-1mm]{\scriptsize stable}};
                  \draw[very thick,dotted] (-4.7,0)--(5.0,0)
                  node[right,align=left] {$0$\\[-1mm]{\scriptsize unstable}};

                  % Representative nonconstant trajectories.
                  \draw[thick,domain=-4.5:4.8,samples=120,smooth]
                  plot (\x,{3.14/(1+exp(-0.85*\x))});
                  \draw[thick,domain=-4.5:4.8,samples=120,smooth]
                  plot (\x,{-3.14/(1+exp(-0.85*\x))});
                  \node at (3.2,1.55) {$0<x_0<\pi$};
                  \node at (3.2,-1.55) {$-\pi<x_0<0$};
                \end{tikzpicture}
            \]
            The drawn curves are qualitative sketches; they indicate the correct
            monotonicity and limiting equilibria, not an explicit formula for the
            solutions.

            For a nonconstant solution, separation of variables yields
            \[
                \frac{e^x}{\sin x}\,dx=dt.
            \]
            Hence, as long as $x(t)$ remains in the interval between the neighboring
            fixed points containing $x_0$,
            \[
                \boxed{
                    t=\int_{x_0}^{x(t)}\frac{e^u}{\sin u}\,du.
                }
            \]
            The antiderivative $\int e^x\csc x\,dx$ is not an elementary function, so
            there is no elementary closed-form expression for $x(t)$. If
            \[
                F(x)=\int^x\frac{e^u}{\sin u}\,du
            \]
            is defined on the interval containing $x_0$, then the solution may be
            written implicitly, or in inverse-function form, as
            \[
                F(x(t))-F(x_0)=t,
                \qquad
                \boxed{x(t)=F^{-1}\bigl(t+F(x_0)\bigr).}
            \]
            Together with the constant solutions $x(t)\equiv n\pi$, this describes all
            solutions.
        \end{solution}

        \noindent\textbf{2.2.7} Analyze the following equation graphically.
        Sketch the vector field on the real line find all the fixed points, classify their
        stability, and sketch the graph of $x(t)$ for different initial conditions.
        Then try to find the analytical solution for $x(t)$, or state that it's impossible.
        The equation in question is:
        \[
            \dot{x} = e^{x}-\cos(x)
        \]
        (Hint: Sketch the graphs of $e^x$ and $\cos(x)$ on the same axes, and look for
        intersections. You won't be able to find the fixed points explicitly, but you still can
        find the qualitative behaviour.)

        \begin{solution}
        \end{solution}

        \noindent\textbf{2.2.8} (Working backwards, from flows to equations)
        Given an equation
        \[
            \dot{x}=f(x),
        \]
        we know how to sketch the corresponding flow on the real line.
        Here you are asked to solve the opposite problem: For the phase
        portrait shown in Figure~1, find an equation that is consistent with
        it. (There are an infinite number of correct answers---and wrong ones
        too.)

        \begin{figure}[h]
            \centering
            \begin{tikzpicture}[scale=1.25, >=stealth]

                % Phase line
                \draw (-3,0) -- (4,0);

                % Flow arrows
                \draw[->, thick] (-2.8,0) -- (-2.0,0);
                \draw[->, thick] (-0.85,0) -- (-0.15,0);
                \draw[->, thick] (1.35,0) -- (0.65,0);
                \draw[->, thick] (2.25,0) -- (3.05,0);

                % Half-stable fixed point at x = -1
                \draw[fill=white, thick] (-1,0) circle (2.5pt);
                \begin{scope}
                    \clip (-1.1,-0.1) rectangle (-1,0.1);
                    \fill (-1,0) circle (2.5pt);
                \end{scope}
                \draw[thick] (-1,0) circle (2.5pt);

                % Stable fixed point at x = 0
                \fill (0,0) circle (2.5pt);

                % Unstable fixed point at x = 2
                \draw[fill=white, thick] (2,0) circle (2.5pt);

                % Labels
                \node[below=5pt] at (-1,0) {$-1$};
                \node[below=5pt] at (0,0) {$0$};
                \node[below=5pt] at (2,0) {$2$};

            \end{tikzpicture}
            \caption{}
        \end{figure}

        \begin{solution}
        \end{solution}

        \noindent\textbf{2.3.2} (Autocatalysis)
        Consider the model chemical reaction
        \[
            A+X
            \underset{k_{-1}}{\overset{k_1}{\rightleftarrows}}
            2X,
        \]
        in which one molecule of \(X\) combines with one molecule of \(A\)
        to form two molecules of \(X\). This means that the chemical \(X\)
        stimulates its own production, a process called \emph{autocatalysis}.
        This positive feedback process leads to a chain reaction, which
        eventually is limited by a ``back reaction'' in which \(2X\) returns
        to \(A+X\).

        According to the \emph{law of mass action} of chemical kinetics, the
        rate of an elementary reaction is proportional to the product of the
        concentrations of the reactants. We denote the concentrations by
        lowercase letters \(x=[X]\) and \(a=[A]\). Assume that there is an
        enormous surplus of chemical \(A\), so that its concentration \(a\)
        can be regarded as constant. Then the equation for the kinetics of
        \(x\) is
        \[
            \dot{x}=k_1ax-k_{-1}x^2,
        \]
        where \(k_1\) and \(k_{-1}\) are positive parameters called
        \emph{rate constants}.

        \begin{enumerate}[label=\alph*)]
            \item Find all the fixed points of this equation and classify
            their stability.
            \item Sketch the graph of \(x(t)\) for various initial values
            \(x_0\).
        \end{enumerate}

        \begin{solution}
        \end{solution}

        \noindent\textbf{2.3.3} (Tumor growth)
        The growth of cancerous tumors can be modeled by the Gompertz law
        \[
            \dot{N}=-aN\ln(bN),
        \]
        where \(N(t)\) is proportional to the number of cells in the tumor,
        and \(a,b>0\) are parameters.

        \begin{enumerate}[label=\alph*)]
            \item Interpret \(a\) and \(b\) biologically.
            \item Sketch the vector field and graph \(N(t)\) for various
            initial values.
        \end{enumerate}

        The predictions of this simple model agree surprisingly well with
        data on tumor growth, as long as \(N\) is not too small; see
        Aroesty et al.\ (1973) and Newton (1980) for examples.

        \begin{solution}
        \end{solution}

        \noindent\textbf{2.3.5} (Dominance of the fittest) Suppose $X$ and $Y$ are two species
        that reproduce exponentially fast: $\dot{X} = aX$ and $\dot{Y} = bY$, respectively, with
        initial conditions $X_0,Y_0 > 0$ and growth rates $a\ge b\ge 0$.
        Here $X$ is ``fitter'' than $Y$ in the sense that it reproduces faster, as reflected by
        the inequality $a > b$.
        So we'd expect $X$ to keep increasing its share of the total population $X + Y$ as
        $t\rightarrow\infty$.
        The goal of this exercise is to demonstrate this intuitive result, first analytically and
        then geometrically.
        \begin{itemize}
            \item[(a)] Let $x(t) = X(t)/[X(t) + Y(t)]$ denote $X$'s share of the total population.
            By solving for $X(t)$ and $Y(t)$, show that $x(t)$ increasing monotonically approaches
            1 as $t\rightarrow\infty$.
            \item[(b)] Alternatively, we can arrive at the same conclusions by deriving a differential
            equation for $x(t)$.
            To do so, take the time derivative of $x(t) = X(t)/[X(t) + Y(t)]$ using the quotient and
            chain rules.
            Then substitute for $X$ and $Y$ and therby show that $x(t)$ obeys the logistic equation
            $\dot{x} = (a - b)x(1 - x)$.
            Explain why this implies that $x(t)$ increases monotonically and approaches 1 as
            $t\rightarrow\infty$.
        \end{itemize}

        \begin{solution}
        \end{solution}

        \noindent\textbf{2.4.3} Use linear stability analysis to classify the fixed points
        of the following system.
        If linear stability analysis fails because $f'(x^*)=0$, use a graphical argument to decide
        the stability.
        \[
            \dot{x} = \tan(x)
        \]

        \begin{solution}
        \end{solution}

        \noindent\textbf{2.4.8} using linear stability analysis, classify the fixed points of the
        Gompertz model of tumor growth $\dot{N} = -aN\ln(bN)$.
        (Similar to Exercise 2.3.3, $N(t)$ is proportional to the number of cells in the tumor and
        $a,b > 0$ are parameters.)

        \begin{solution}
        \end{solution}

        \noindent\textbf{2.8.2} Sketch the slope field for the following differential equations.
        Then ``integrate'' the equation manually by drawing trajectories that are everywhere
        parallel to the local slope.
        \begin{itemize}
            \item[(a)] $\dot{x} = x$
            \item[(c)] $\dot{x} = 1 - 4x(1 - x)$
        \end{itemize}

        \begin{solution}
        \end{solution}

        \noindent\textbf{2.8.6} (Analytically intractible problem) Consider the initial value
        problem $\dot{x} = x + e^{-x}$, $x(0) = 0$.
        \begin{itemize}
            \item[(a)] Sketch the solution $x(t)$ for $t\ge 0$
            \item[(b)] Using some analytical arguments, obtain rigorous bounds on the value of $x$ at
            $t=1$.
            In other words profe that $a < x(1) < b$, for $a,b$ to be determined.
            By being clever, try to make $a$ and $b$ as close together as possible.
            (Hint: Bound the given vector field by approximate vector fields that can be integrated
            analytically.)
            \item[(c)] Now for the numerical part: Using the Euler method, compute $x$ at $t=1$, correct
            to three decimal places.
            How small does the stepsize need to be to obtain the desired accuracy?
            (Give the order of magnitude, not the exact number)
        \end{itemize}

        \begin{solution}
        \end{solution}

    \end{onehalfspace}
\end{document}